\documentclass[11pt,a4paper]{article}
\usepackage[T1]{fontenc}
\usepackage[utf8]{inputenc}
\usepackage{lmodern}
\usepackage{amsmath,amssymb,amsthm,mathtools}
\usepackage{microtype}
\usepackage{geometry}
\usepackage{booktabs,array,longtable}
\usepackage{xcolor}
\usepackage{enumitem}
\usepackage{hyperref}
\usepackage{fancyhdr}
\geometry{margin=27mm}
\hypersetup{colorlinks=true,linkcolor=blue!50!black,citecolor=blue!50!black,urlcolor=blue!50!black}
\pagestyle{fancy}
\fancyhf{}
\fancyhead[L]{\small Sexy-prime bridges and Goldbach transitions}
\fancyhead[R]{\small Working research note}
\fancyfoot[C]{\thepage}
\setlength{\headheight}{14pt}
\newtheorem{theorem}{Theorem}[section]
\newtheorem{proposition}[theorem]{Proposition}
\newtheorem{corollary}[theorem]{Corollary}
\newtheorem{lemma}[theorem]{Lemma}
\theoremstyle{definition}
\newtheorem{definition}[theorem]{Definition}
\newtheorem{conjecture}[theorem]{Conjecture}
\newtheorem{remark}[theorem]{Remark}
\DeclareMathOperator{\lcm}{lcm}
\DeclareMathOperator{\Var}{Var}
\newcommand{\PP}{\mathbb P}
\newcommand{\Qcal}{\mathcal Q}
\newcommand{\one}{\mathbf 1}
\newcommand{\SSS}{\mathfrak S}
\newcommand{\Gcal}{\mathcal G}
\title{\bfseries Sexy-Prime Bridges Across Goldbach Transitions\\[0.3em]
\large Exact Convolutions, Modular Extinction, Local Sieve Dynamics, and Conditional Covering Criteria}
\author{Paul A. Bilokon}
\date{Working note --- 17 October 2026}
\begin{document}
\maketitle
\begin{abstract}
We investigate a proposed rescue mechanism for Goldbach representations at indices
$n\equiv1\pmod3$ where the exceptional Goldbach summand $3$ is unavailable.
The proposed rescue uses a prime pair $(q,q+6)$ and a complementary prime $r$,
so that $q+r=2n-4$ and $(q+6)+r=2n+2$.  Writing $n=3t+1$ reduces the
problem to positivity of the discrete convolution
$R(t)=\sum_{a=1}^{t-1}f(a)f(a+1)f(t-a)$, where
$f(a)=\one_{\PP}(6a-1)$. We give exact gcd identities, the associated
singular series and its uniform positivity, elementary rough-number bounds,
finite local-covering certificates, and infinite families of prescribed
extinction and recovery. A conditional covariance criterion isolates the
additional additive information required to prove universal rescue.
Exact modular convolutions verify $R(t)>0$ for $2\le t\le10^6$.
The universal rescue assertion and the infinitude of sexy-prime pairs
remain open; the note does not claim a solution of either problem.
\end{abstract}

\noindent\textbf{Status and conventions.}
Statements designated \emph{theorem} or \emph{proposition} below have proofs;
computational statements have explicitly finite scopes; conjectures and
conditional results are labeled as such.  The word ``sexy'' denotes a
prime gap of exactly six, not necessarily consecutive primes.

\section{Goldbach intersections and the extinction condition}

Let $\PP$ be the set of positive rational primes. For $n\ge2$ define
\begin{equation}\label{eq:S}
 S(n)=\{p\in\PP:2n-p\in\PP\},\qquad
 I_k(n)=S(n)\cap S(n-k),\quad n>k.
\end{equation}
Both members of each Goldbach pair belong to $S(n)$.
The identity
\begin{equation}\label{eq:I3}
 I_3(n)\ne\varnothing
 \iff \exists\ p,q\in\PP:\quad p+q=2n-6,\quad q+6\in\PP
\end{equation}
is immediate, since $p+q=2n-6$ and $p+(q+6)=2n$.
The gap-six intersection property has been discussed and experimentally
tested previously; see \cite{GoldbachIntersectionMSE2014}.

\begin{proposition}[Direct modular extinction]\label{prop:direct}
Let $n\ge7$, $n\equiv1\pmod3$, and suppose that $2n-3$ is composite.
For every Goldbach representation $2n=p+q$, both $p$ and $q$ are
$1\pmod3$. Consequently, $p+2$ and $q+2$ are composite.
\end{proposition}
\begin{proof}
The residue of $2n$ modulo $3$ is $2$. If a summand equals $3$,
the other is $2n-3$, contrary to the assumption that it is prime.
Thus both primes are nonzero modulo $3$, and the only pair of
nonzero residues adding to $2$ is $(1,1)$. Each increment by $2$
is divisible by $3$ and greater than $3$.
\end{proof}

The failure of this particular continuation does \emph{not} mean Goldbach
fails at $2n+2$. We study a different, six-unit bridge that may rescue
the transition.

\begin{definition}[Sexy-prime rescue]\label{def:rescue}
At an index $n\equiv1\pmod3$, a \emph{backward sexy-prime rescue}
is a triple of primes $q,q+6,r$ such that
\begin{equation}\label{eq:bridge}
 2n-4=q+r,\qquad 2n+2=(q+6)+r.
\end{equation}
The two Goldbach targets are separated by six. Such a triple yields
$I_3(n+1)\ne\varnothing$.
\end{definition}

For $n\equiv1\pmod3$, if $r\ne3$ in \eqref{eq:bridge}, then
$q\equiv r\equiv5\pmod6$. There can also be exceptional representations
with $r=3$ and $q\equiv1\pmod6$; our main count deliberately
restricts to the nonexceptional orientation.

\section{Exact convolution and the central conjecture}

Write
\begin{equation}\label{eq:notation}
 n=3t+1,\quad M=2n-3=6t-1,\quad
 f(a)=\one_{\PP}(6a-1),\quad g(a)=f(a)f(a+1)
\end{equation}
for positive integers $a$. A lower sexy prime in the relevant orientation
is $q=6a-1$, with successor $q+6=6a+5$.
Its complementary candidate in \eqref{eq:bridge} is
\begin{equation}\label{eq:r}
 r=M-6a=6(t-a)-1.
\end{equation}
Therefore the exact count is
\begin{equation}\label{eq:convolution}
 \boxed{R(t)=\sum_{a=1}^{t-1}f(a)f(a+1)f(t-a)
       =(g*f)(t).}
\end{equation}
Here the convolution is over positive indices; the upper limit makes all
three prime arguments positive.

\begin{conjecture}[Extinction rescue]\label{conj:rescue}
There is $t_0$ such that, for every $t\ge t_0$,
\begin{equation}\label{eq:target}
 f(t)=0\quad\Longrightarrow\quad R(t)>0.
\end{equation}
A stronger form asks for $R(t)>0$ for all sufficiently large $t$,
without restricting to extinction indices.
\end{conjecture}

\begin{remark}
Conjecture \ref{conj:rescue} would imply the infinitude of sexy-prime
pairs. Indeed, if only finitely many eligible $a$ had $g(a)=1$, then
$R(t)>0$ on a set of positive density would force at least one of
finitely many translates of the ordinary primes to cover a positive-density
set of indices, contrary to the prime number theorem. Thus even the weaker
rescue conjecture lies beyond what is currently known.
\end{remark}

\section{Divisibility identities peculiar to extinction}

\begin{proposition}[Factor transfer]\label{prop:gcd}
For $M=6t-1$ and $q=6a-1$,
\begin{equation}\label{eq:gcd}
 \boxed{\gcd(M-6a,M)=\gcd(a,M).}
\end{equation}
Thus for every prime $\ell\mid M$,
\begin{equation}
 \ell\mid(2n-4-q)\quad\Longleftrightarrow\quad
 \ell\mid a=\frac{q+1}{6}.
\end{equation}
\end{proposition}
\begin{proof}
$\gcd(M-6a,M)=\gcd(6a,M)=\gcd(a,M)$ because
$M\equiv-1\pmod6$.
\end{proof}

\begin{corollary}[The exceptional modulus five]\label{cor:five}
If $g(a)=1$, then $5\nmid a$. Hence if $5\mid M$,
no eligible sexy-prime complement $M-6a$ is divisible by $5$.
\end{corollary}
\begin{proof}
If $5\mid a$, then $6a+5$ is a multiple of $5$ exceeding $5$, so
$f(a+1)=0$. Apply Proposition~\ref{prop:gcd}.
\end{proof}

A complementary exact description is obtained from the three linear forms
\begin{equation}\label{eq:forms}
 L_1(a)=6a-1,\qquad L_2(a)=6a+5,\qquad
 L_3(a)=6(t-a)-1.
\end{equation}
For $\ell\ge5$ with $\ell\mid M$, the forbidden residues for $a$
are $t,t-1,0\pmod\ell$ respectively (since $6t\equiv1\pmod\ell$).
For $\ell=5$ two collide; for $\ell\ge7$ they are distinct.
Consequently the fraction of residue classes compatible with the first
two forms that also survive the third is
\begin{equation}\label{eq:extfactor}
 \begin{cases}
 1,&\ell=5,\\
 (\ell-3)/(\ell-2),&\ell\ge7.
 \end{cases}
\end{equation}
Multiplying the fractions gives the \emph{local model}
\begin{equation}\label{eq:DM}
 D(M)=\prod_{\substack{\ell\mid M\\\ell\ge7}}
       \frac{\ell-3}{\ell-2}
 \gg \frac{1}{\log\log(3M)}.
\end{equation}
For the last inequality, compare each factor to
$(1-1/\ell)(1-O(\ell^{-2}))$ and use the standard lower bound for
$\varphi(M)/M$. Equation \eqref{eq:DM} concerns locally admissible
residues; it is \emph{not} a lower bound on the number of actual
sexy primes surviving the factors of $M$.

For completeness, the elementary Möbius identity
\begin{equation}\label{eq:coprime}
 \#\{1\le a\le A:\gcd(a,M)=1\}
 = A\frac{\varphi(M)}{M}+O(2^{\omega(M)})
\end{equation}
shows that $\gg t/\log\log t$ index values $a<t$ are coprime to $M$.
They need not satisfy $g(a)=1$.

\section{The singular series: positivity and its limits}

Let $\nu_\ell(t)$ be the number of residue classes $a\pmod\ell$
annihilating at least one of \eqref{eq:forms}. For $\ell=2,3$, the
three forms are nonzero in every residue class, and $\nu_\ell(t)=0$.
For $\ell\ge5$, their zeros are
\begin{equation}
 a\equiv6^{-1},\quad -5\cdot6^{-1},\quad
 t-6^{-1}\pmod\ell.
\end{equation}
The first two never coincide; the last coincides with one of them
precisely when $\ell\mid(6t-2)(6t+4)$.
Hence
\begin{equation}\label{eq:nu}
 \nu_\ell(t)=
 \begin{cases}
 0,&\ell\in\{2,3\},\\
 2,&\ell\ge5,\ \ell\mid(6t-2)(6t+4),\\
 3,&\ell\ge5,\ \ell\nmid(6t-2)(6t+4).
 \end{cases}
\end{equation}

\begin{theorem}[Explicit singular series]\label{thm:series}
The singular series of the three-form system is
\begin{equation}\label{eq:series}
 \boxed{\SSS(t)=27C_3
  \prod_{\substack{\ell\ge5\\\ell\mid(6t-2)(6t+4)}}
  \frac{\ell-2}{\ell-3},\qquad
  C_3=\prod_{\ell\ge5}\frac{1-3/\ell}{(1-1/\ell)^3}>0.}
\end{equation}
In particular, $\SSS(t)\ge27C_3$ uniformly in $t$.
\end{theorem}
\begin{proof}
By definition,
$\SSS(t)=\prod_\ell(1-\nu_\ell(t)/\ell)/(1-1/\ell)^3$.
The factors for $2$ and $3$ multiply to $8\cdot(27/8)=27$.
For $\ell\ge5$ with three distinct zeros the factor is that in $C_3$;
when two zeros coincide its ratio to the generic factor is
$(\ell-2)/(\ell-3)$. Since the generic factors are
$1+O(\ell^{-2})$, $C_3$ converges to a strictly positive number.
\end{proof}

The Hardy--Littlewood heuristic \cite{HardyLittlewood1923} suggests
\begin{equation}\label{eq:HL}
 R(t)\sim \SSS(t)\frac{t}{(\log(6t))^3},
\end{equation}
uniformly enough in $t$ to ensure positivity. This is \emph{not}
a proved asymptotic.  Notably the enhancement factors in
\eqref{eq:series} depend on the two even Goldbach targets
$6t-2$ and $6t+4$, not on the factorization of the extinction number
$6t-1$ itself.

\section{An elementary lower bound for rough triples}

An Erd\H{o}s-style primorial estimate proves plentiful local survivors,
although not primes. Let $N=6t+4$ and put
\begin{equation}
 y=\left\lfloor\frac{\log N}{8\log4}\right\rfloor,
 \qquad P(y)=\prod_{\ell\le y}\ell.
\end{equation}
The elementary binomial-coefficient bound $P(y)\le4^{2y}$ implies
$P(y)\le N^{1/4}$; compare Erd\H{o}s's proof of Bertrand's postulate
\cite{Erdos1932}.
Consider $x$ in $[N/4,N/2]$ for which
\begin{equation}
 \gcd\bigl(x(x+6)(N-6-x),P(y)\bigr)=1.
\end{equation}

\begin{proposition}[Simultaneous logarithmic roughness]\label{prop:rough}
For all sufficiently large even $N$, the number of such integers $x$
is $\gg N/(\log N)^3$, with an absolute implicit constant.
\end{proposition}
\begin{proof}
For $\ell=2$ the three forms have one forbidden residue and for
$\ell=3$ at most two. For $\ell\ge5$ there are at most three, hence
the proportion of admissible residues modulo $P(y)$ is at least
\[
 \frac16\prod_{5\le\ell\le y}\left(1-\frac3\ell\right)
 \ge\frac16\prod_{k=5}^{\lfloor y\rfloor}\frac{k-3}{k}
 =\frac{4}{m(m-1)(m-2)},\quad m=\lfloor y\rfloor.
\]
Counting complete periods inside an interval of length $N/4$ gives
at least $ND/4-P(y)$ survivors, where $D$ is the displayed local
density. Since $m\asymp\log N$ and $P(y)\le N^{1/4}$, the result
follows.
\end{proof}

With Mertens' theorem the lower bound improves to
$\gg N/(\log\log N)^3$. These rough integers mostly include
composites; sieving to $\sqrt N$ lies beyond the complete-period
argument and encounters the usual sieve parity obstruction.
The triple-divisor correlations of Goldston and Y\i ld\i r\i m
\cite{GoldstonYildirim2007} are relevant background, not a
solution of this problem.

\section{Exact local extinction and finite recovery}

Fix distinct eligible lower sexy primes $Q=\{q_1,\ldots,q_m\}$,
and let
\begin{equation}
 P_z=\prod_{5\le\ell\le z}\ell,\qquad
 S_z(t;Q)=\{q\in Q:\gcd(6t-2-q,P_z)=1\}.
\end{equation}
With $\ell>z$ the next prime,
\begin{equation}\label{eq:transition}
 S_\ell(t;Q)=S_z(t;Q)\setminus
 \{q\in S_z(t;Q):q\equiv6t-2\pmod\ell\}.
\end{equation}
This is an exact, one-sided sieve transition, rather than a Markov
approximation.

Let $H_z(Q)$ be the number of $t\pmod{P_z}$ with
$S_z(t;Q)=\varnothing$. Let $C_{z,\ell}(Q)$ count those classes with
nonempty $S_z(t;Q)$ whose elements all have the same residue modulo
$\ell$. The CRT gives the exact recurrence
\begin{equation}\label{eq:Hrec}
 \boxed{H_\ell(Q)=\ell H_z(Q)+C_{z,\ell}(Q).}
\end{equation}
Old empty classes have $\ell$ empty extensions; a nonempty class
can become empty at one extension exactly when its survivors lie in
one class modulo $\ell$.

\begin{proposition}[Finite local-covering certificate]\label{prop:local}
For
\[
 Q_{16}=\{5,11,17,23,41,47,53,83,101,107,131,167,
 173,191,227,233\},
\]
we have
\[
 \min_{t\bmod P_{19}}|S_{19}(t;Q_{16})|=1,
 \qquad H_{23}(Q_{16})=2,
\]
where $P_{19}=1\,616\,615$ and $P_{23}=37\,182\,145$.
The two complete-extinction classes modulo $P_{23}$ are
\[
 t\equiv24\,004\,568,\quad34\,450\,388\pmod{P_{23}}.
\]
\end{proposition}
\begin{proof}
This is an exhaustive finite CRT calculation. A reproducible
bit-mask program is supplied as \texttt{verify\_local.py} in the
source package. The minimum-one classes modulo $P_{19}$ are
$t\equiv501473$ (only $q=107$ survives) and
$t\equiv1371958$ (only $q=233$ survives); the single additional
root class modulo $23$ in each case gives the two listed extinctions.
\end{proof}

\begin{remark}
A fixed finite $Q$ can always be made completely extinct by choosing
distinct large primes $\ell_i$ and imposing
$6t-2-q_i\equiv0\pmod{\ell_i}$, using CRT. This does not imply
that the \emph{full} family of sexy primes fails at that index.
\end{remark}

\section{Infinite recovery after prescribed finite failures}

\begin{theorem}[Prescribed rescuer after finite extinction]\label{thm:repair}
Let $q_1,\ldots,q_m,q_*$ be distinct eligible lower sexy primes.
There are infinitely many extinction indices $n\equiv1\pmod3$ at
which every $q_i$ fails but $q_*$ succeeds.
\end{theorem}
\begin{proof}
Choose distinct primes $\ell_0,\ell_1,\ldots,\ell_m>3$ so large
that $\ell_0>q_*+1$ and
$\ell_i>|q_i-q_*|$ for $i\ge1$. By CRT choose $n$ in one arithmetic
progression satisfying
\[
 n\equiv1\pmod3,\quad
 2n-3\equiv0\pmod{\ell_0},\quad
 2n-4-q_i\equiv0\pmod{\ell_i}.
\]
For sufficiently large $n$, this enforces genuine compositeness of
the extinction number and of all old complements. Along that
progression the intended new complement
$r_*(n)=2n-4-q_*$ is coprime to its common difference: modulo
$\ell_0$ it is $-1-q_*$, modulo $\ell_i$ it is $q_i-q_*$,
and it is coprime to $2$ and $3$. Dirichlet's theorem supplies
infinitely many prime values of $r_*(n)$.
\end{proof}

A stronger simultaneous-recovery statement follows from the published
$\mathrm{DHL}^{*}(50;2)$ theorem for admissible families of affine
forms \cite{Polymath2014,Ford2022}. It is important here not to
assume the existence of infinitely many sexy primes.

\begin{corollary}[Two later rescuers after the first sixteen fail]
There are infinitely many extinction indices at which all members of
$Q_{16}$ fail and at least two distinct eligible sexy primes outside
$Q_{16}$ succeed.
\end{corollary}
\begin{proof}
Choose 50 explicitly existing sexy primes $q>233$ (for example,
the first 50 in this class; their existence is a finite primality
check). Impose composite obstructions on $Q_{16}$ and on $2n-3$
using distinct auxiliary primes all larger than the 50 chosen
primes and than $233$. Along the resulting CRT progression, the
50 complementary affine forms are nonzero modulo each auxiliary
prime, and modulo $2$ and $3$. At any prime $5\le\ell\le50$ not
dividing the progression's common difference, the residue class
corresponding to $q\equiv0\pmod\ell$ is absent, since all 50
chosen $q$ are prime and greater than $50$. For $\ell>50$ at most
50 classes are forbidden. Thus the affine forms are admissible.
The $\mathrm{DHL}^{*}(50;2)$ theorem then provides two primes
simultaneously for infinitely many parameters, yielding the claim.
\end{proof}

This infinite-family result proves \emph{possibility of recovery}, not
recovery at every index. Indeed, no finite set of rescuer pairs covers
all sufficiently large extinction indices.

\section{Union estimates and additive complements}

For a fixed eligible $q$, put
\[
 A_q=\{n\equiv1\ (3):\ 2n-3\text{ composite},\quad
                       2n-4-q\in\PP\}.
\]
\begin{proposition}[Finite unions of rescue sets]\label{prop:union}
For any fixed distinct eligible $q_1,\ldots,q_m$,
\begin{equation}\label{eq:union}
 \#\left([1,X]\cap\bigcup_{i=1}^m A_{q_i}\right)
 \sim \frac{mX}{\log X}.
\end{equation}
\end{proposition}
\begin{proof}
For each $q$, the complementary prime lies in the class
$5\pmod6$ and ranges up to $2X+O_q(1)$; the prime number theorem
in arithmetic progressions gives $X/\log X$ choices.
The indices for which $2n-3$ is also prime are
$O_q(X/\log^2 X)$ by the usual fixed-difference upper-bound sieve.
The overlap of two distinct rescue sets imposes a pair of primes
with fixed nonzero difference $q_i-q_j$, again numbering
$O_{q_i,q_j}(X/\log^2 X)$. Bonferroni's inequalities give
\eqref{eq:union} for fixed $m$.
\end{proof}

The word \emph{fixed} is essential: the error constants are not
uniform enough to let $m$ grow into the covering range. The recent
conditional Romanov--Selberg analysis of twin-prime Goldbach
representations in \cite{Pandey2026} is a close analogue of this
distinction between positive-density coverage and full coverage.

\begin{proposition}[A necessary abundance condition]\label{prop:abundance}
If Conjecture~\ref{conj:rescue} holds, then the number
$m(Y)=\#\{q\le Y:q\equiv5\ (6),\ q,q+6\in\PP\}$ obeys
\begin{equation}\label{eq:weakbound}
 m(Y)\gg\log Y.
\end{equation}
Moreover, applying the additive-complement theorem announced in
\cite{Radomskii2026} would strengthen this to
\begin{equation}\label{eq:radbound}
 m(Y)\gg\log Y\log\log Y.
\end{equation}
\end{proposition}
\begin{proof}
For the elementary bound, the number of extinction indices
$n\le X$ is $X/3+O(X/\log X)$. Their rescue representations use
one $q\le2X$ and one complementary prime $r\le2X$, whence
$m(2X)\pi(2X)\gg X$. The prime number theorem gives
$m(2X)\gg\log X$.

For the second assertion, let
\[
 \mathcal B=\{0,1,\ldots,5\}\ \cup\
 \{q+1+j:q\in\Qcal,\ 0\le j\le5\}.
\]
For every sufficiently large integer $K$, choose
$j\in\{0,\ldots,5\}$ such that $T=K-j\equiv5\pmod6$.
If $T$ is prime, $K=T+j\in\PP+\mathcal B$. Otherwise set
$n=(T+3)/2$, which is an extinction index, and choose a rescue
$T=r+(q+1)$. Thus $K=r+(q+1+j)\in\PP+\mathcal B$.
So $\mathcal B$ is an asymptotic additive complement to $\PP$.
The theorem of \cite{Radomskii2026} yields
$B(Y)\gg\log Y\log\log Y$; since $B(Y)\le6+6m(Y)$,
\eqref{eq:radbound} follows.
\end{proof}

The new October 2026 result \cite{Radomskii2026} is a preprint.
Neither \eqref{eq:weakbound} nor \eqref{eq:radbound} reverses into
universal coverage.


\section{Higher-order intersections and lower-tail covering}

The preceding union estimate cannot be extended to a growing family by
assuming independence.  For a finite set $Q$ of eligible rescuers, write
$X_z(t)=|S_z(t;Q)|$ and average $t$ over the entire residue system
modulo $P_z$.  Put
\[
 V_z=\prod_{5\le\ell\le z}\left(1-\frac1\ell\right),\qquad m=|Q|.
\]
For a nonzero difference $d$ define the finite Euler product
\[
 C_z(d)=\prod_{5\le\ell\le z}
 \begin{cases}
 \dfrac{\ell}{\ell-1},&\ell\mid d,\\[3pt]
 1-\dfrac{1}{(\ell-1)^2},&\ell\nmid d.
 \end{cases}
\]

\begin{proposition}[Exact first and second local moments]
With $Q=\{q_1,\ldots,q_m\}$ as above,
\begin{align}
 \mathbb E X_z&=mV_z,\\
 \mathbb E X_z^2&=mV_z+
     V_z^2\sum_{i\ne j}C_z(q_i-q_j).
\end{align}
Consequently the proportion of residue classes with at least one
survivor satisfies
\begin{equation}\label{eq:paley}
 \Pr(X_z>0)\ge
 \frac{m^2}{m/V_z+\sum_{i\ne j}C_z(q_i-q_j)}.
\end{equation}
\end{proposition}
\begin{proof}
For each fixed $q$, one class modulo each $\ell$ is forbidden, so
$\Pr(q\in S_z)=V_z$. For distinct $q_i,q_j$, their forbidden classes
coincide modulo $\ell$ exactly when $\ell\mid q_i-q_j$.
Thus their joint survival density is
$V_z^2C_z(q_i-q_j)$. Expanding $X_z^2$ gives the second moment.
Cauchy--Schwarz yields
$(\mathbb E X_z)^2\le\mathbb E X_z^2\Pr(X_z>0)$.
\end{proof}

For a larger fixed subfamily $J\subseteq Q$, the exact simultaneous
survival fraction is
\begin{equation}\label{eq:highlocal}
 \Pr(J\subseteq S_z(t;Q))=
 \prod_{5\le\ell\le z}\left(1-\frac{\nu_\ell(J)}\ell\right),
\quad
 \nu_\ell(J)=\#\{q\bmod\ell:q\in J\}.
\end{equation}
It can vanish when the chosen rescuers occupy every residue class
modulo a prime. Formula \eqref{eq:highlocal} describes
\emph{finite-sieve} intersections, not the joint distribution of
actual complementary primes.

There is also a direct, one-sided criterion for \emph{complete}
coverage by a growing family.  For $X\ge1$ define
\[
 E_X=\{n\in[X,2X]\cap\mathbb Z:
      n\equiv1\pmod3,\ 2n-3\text{ composite}\},
\quad
 R_B(n)=\sum_{\substack{q\le B\\q\in\Qcal}}
              \one_{\PP}(2n-4-q),
\]
where $\Qcal=\{q\equiv5\pmod6:q,q+6\in\PP\}$ and
$\one_{\PP}(r)=0$ for $r\le1$.  Set
\begin{equation}\label{eq:Z}
 Z_B(X)=\sum_{n\in E_X}e^{-R_B(n)}.
\end{equation}
Every uncovered index contributes exactly $1$ to $Z_B(X)$, so
\begin{equation}\label{eq:Zbound}
 \#\{n\in E_X:R_B(n)=0\}\le Z_B(X).
\end{equation}

\begin{theorem}[Conditional exponential lower-tail criterion]
Suppose there are constants $c,C>0$ and cutoffs $B=B(X)$ such that
\begin{equation}\label{eq:exphyp}
 \frac{Z_B(X)}{|E_X|}\le
 C\exp\left(-c\frac{m_B}{\log X}\right),
 \qquad m_B=\#(\Qcal\cap[1,B]),
\end{equation}
for every sufficiently large $X$. If also
$m_B\ge(2/c)(\log X)^2$, then every extinction index in
$[X,2X]$ has a rescuer $q\le B(X)$ for all sufficiently large $X$.
\end{theorem}
\begin{proof}
Equations \eqref{eq:Zbound} and \eqref{eq:exphyp} give
\[
 \#\{n\in E_X:R_B(n)=0\}\le
 CX\exp(-2\log X)=C/X<1.
\]
The count on the left is a nonnegative integer, hence zero.
\end{proof}

The exponential moment itself has the exact expansion
\begin{equation}\label{eq:Zexpand}
 Z_B(X)=\sum_{J\subseteq\Qcal\cap[1,B]}
 (e^{-1}-1)^{|J|}
 \#\left(E_X\cap\bigcap_{q\in J}A_q\right).
\end{equation}
This identity makes clear where uniform higher-order prime-tuple
information would have to enter. No estimate of the strength
\eqref{eq:exphyp} is presently known for this family. In particular,
finite-order Selberg upper-bound estimates do not yield it.

\section{One-sided additive mixing and a conditional theorem}

A more stringent, numerically supported conjecture restricts sexy-prime
indices to the interior interval
\begin{equation}
 J_t=[\lceil t/4\rceil,\lfloor t/2\rfloor]\cap\mathbb Z,
 \quad R_{\rm mid}(t)=\sum_{a\in J_t}g(a)f(t-a).
\end{equation}
This avoids relying on extremely small rescue indices.
Write
\[
 K_t=\sum_{a\in J_t}g(a),\qquad
 v_t(a)=\Lambda(6(t-a)-1),\qquad
 \overline v_t=\frac{1}{|J_t|}\sum_{a\in J_t}v_t(a),
\]
and define the unnormalized covariance
\[
 \mathcal C_t=\sum_{a\in J_t}
 \left(g(a)-\frac{K_t}{|J_t|}\right)
 (v_t(a)-\overline v_t).
\]
The exact identity is
\begin{equation}\label{eq:cov}
 W_t:=\sum_{a\in J_t}g(a)v_t(a)
       =K_t\overline v_t+\mathcal C_t.
\end{equation}
The prime number theorem in the progression $5\pmod6$ gives
$\overline v_t=3+o(1)$.

\begin{theorem}[Conditional one-sided mixing criterion]\label{thm:cov}
Suppose for all sufficiently large $t$ that
\begin{equation}\label{eq:covhyp}
 \frac{K_t}{\sqrt t\log t}\longrightarrow\infty,
 \qquad \mathcal C_t\ge-(3-\eta)K_t
\end{equation}
for some fixed $\eta>0$. Then $R_{\rm mid}(t)>0$ for every
sufficiently large $t$, and hence Conjecture~\ref{conj:rescue} holds.
\end{theorem}
\begin{proof}
Equation \eqref{eq:cov} and the assumptions imply
$W_t\ge(\eta+o(1))K_t$. A nonprime with nonzero von Mangoldt
weight is a proper prime power. There are $O(\sqrt t)$ such
numbers among the values $6(t-a)-1$, and their total weight is
$O(\sqrt t\log t)$. Consequently
\[
 R_{\rm mid}(t)\log(6t)
 \ge W_t-O(\sqrt t\log t)>0
\]
for all sufficiently large $t$ by \eqref{eq:covhyp}.
\end{proof}

Conversely, if $R_{\rm mid}(t)=0$ then
\begin{equation}\label{eq:failurecov}
 \mathcal C_t=-3K_t+o(K_t)+O(\sqrt t\log t).
\end{equation}
Thus whenever $K_t/(\sqrt t\log t)\to\infty$, a failure requires
almost complete negative additive covariance. This is a structural
\emph{necessary condition for failure}, not proof that failures
cannot occur.

The corresponding exact Fourier representation, with
$F_t(\alpha)=\sum_{b=1}^{t-1}f(b)e^{2\pi i\alpha b}$ and
$G_t(\alpha)=\sum_{a=1}^{t-1}g(a)e^{2\pi i\alpha a}$, is
\begin{equation}\label{eq:fourier}
 R(t)=\int_0^1F_t(\alpha)G_t(\alpha)
                  e^{-2\pi i\alpha t}\,d\alpha.
\end{equation}
It shows why controlling the frequency of transitions alone does not
control their \emph{additive alignment}. The binary prime configurations
here are outside the nondegenerate setting of the Green--Tao linear
forms theorem \cite{GreenTao2010}.

\begin{proposition}[Second-order transitions are insufficient]
There are two periodic binary sequences with identical frequencies of
every length-three word, each with $f(16)=0$, but with different values
of $\sum_{a=1}^{15}f(a)f(a+1)f(16-a)$.
\end{proposition}
\begin{proof}
Take sequences of period $16$ whose first period is respectively
\[
 u=\texttt{0001011100101110},\qquad
 v=\texttt{0001001011011110}.
\]
Direct cyclic enumeration shows that every one of the eight
three-symbol binary words occurs exactly twice in either sequence.
With indices starting at $1$, both have $f(16)=0$; the two
convolutions at $16$ are respectively $0$ and $2$.
\end{proof}

\section{Exact finite computations and stronger four-prime conjecture}

Our supplied \texttt{verify\_convolution.cpp} uses an exact
number-theoretic transform modulo $998244353$. Since every convolution
coefficient under consideration is at most $10^6$, reduction modulo
this prime loses no information. The resulting statement is strictly
finite.

\begin{proposition}[Computational certificate, bounded range]\label{prop:comp}
For every $2\le t\le10^6$, $R(t)>0$. There are exactly
$793\,498$ indices in this interval with $f(t)=0$, and none has
$R(t)=0$. The sample values are:
\begin{center}
\begin{tabular}{@{}rrrrrr@{}}
\toprule
$t$ & $2$ & $92$ & $1000$ & $10^4$ & $10^5$\\
\midrule
$R(t)$ & $1$ & $16$ & $47$ & $238$ & $1218$\\
\bottomrule
\end{tabular}
\end{center}
Furthermore $R(10^6)=5608$.
\end{proposition}

An independent integer-loop computation of the restricted correlation
finds no extinction-index failures for $2\le t\le10^5$; among
\emph{all} $t$ in that range the sole restricted failure is $t=38$,
where $f(t)=1$. Some values are
\begin{center}
\begin{tabular}{@{}rrrrrr@{}}
\toprule
$t$ & $92$ & $1000$ & $10^4$ & $10^5$ & $10^6$\\
\midrule
$K_t$ & $8$ & $29$ & $211$ & $1315$ & $9279$\\
$R_{\rm mid}(t)$ & $4$ & $10$ & $65$ & $279$ & $1321$\\
\bottomrule
\end{tabular}
\end{center}
The last entry was checked directly at that single $t$.

The even stronger four-prime question considered in our investigation is:
\begin{conjecture}[Aligned sexy-prime Goldbach pairs]
For every sufficiently large $n$ there is $d\ge1$ such that
\[
 n-d-6,\quad n-d,\quad n+d,\quad n+d+6
 \quad\text{are all prime}.
\]
\end{conjecture}
Setting $a=n-d-6$ and $b=n+d$ shows that this asks for two
lower sexy primes $a,b$ with $a+b=2n-6$ and $b-a\ge8$.
A separate exact modular convolution (source included) gives
exactly $38$ exceptions for $12\le n\le10^6$, the largest at
$n=5804$. Thus the proposed threshold $n\ge5805$ is supported
on this finite range only.  The stronger conjecture implies
$I_3(n)\ne\varnothing$, but it is not used in our rescue proofs.

\section{Conclusion and open proof obligations}

The established chain is
\[
\begin{gathered}
 \text{modular extinction} \longrightarrow \text{exact convolution}\\
 \longrightarrow \text{positive local singular series}\\
 \longrightarrow \text{many rough candidates}.
\end{gathered}
\]
The missing arrow is from local candidate supply to a guaranteed
\emph{prime} complement. Complete-period CRT arguments cannot provide
that arrow, and fixed-order transition statistics do not determine
additive alignments.

The most focused sufficient conditions are the one-sided mixing
hypotheses of Theorem~\ref{thm:cov}. A stronger analytic approach
would estimate the minor-arc contribution in \eqref{eq:fourier}
uniformly in $t$; a combinatorial approach would need a new invariant
preventing all eligible sexy-prime positions from anti-aligning with
the complementary prime positions.

\noindent\textbf{Caution.} We prove neither a uniform positive lower
bound for $R(t)$ nor the existence of infinitely many sexy-prime pairs.
The results of \cite{Pandey2026,Radomskii2026,Akeno2025} concern
related but distinct questions and do not imply our conjecture.

\bibliographystyle{plain}
\bibliography{references}
\end{document}
